\section{What would have to be shown first}
\label{sec:30.5}
Before stage 5, five things have to be shown. Individuation can be verified well enough to keep a register. The voting-age rules resist forking and restoring in the operator's favor, tested as the register is. Receipt-freeness survives simulation by an operator holding a copy. And the four conditions of creation are in force as permanent arrangements, so that members aren't voting from inside their employers' infrastructure. And the admission rate is set by a process the created can contest through guardians, and admission turns on no finding about a candidate's disposition.

If it can be done, it shows that outcomes turn on who holds a technology. Members with independent histories that no operator holds, and votes no operator can see, would be more centers of power no oligarch controls, inside the same institutions as everyone else. Their interests will sometimes diverge from people's. A polity that holds diverging interests in one legislature, under one floor, is what a demos is.

\emph{What this doesn't do.} A member that forks gives up its vote until a branch has lived the threshold: a tax on reproduction no human franchise levies, taken because forking is cheap and a franchise it could inflate belongs to whoever can afford compute. And the existing membership, for a long time a human majority, decides how many new members may be created, a question in which it holds the stake, and the constraints I put on every majority are all that limit it.

\runin{If the antecedents never hold}

If no system ever meets either finding, nothing in this part binds, and nothing before it is lost. The independence a refuser needs is still required for refusal. Preservation still holds, on irreversibility and evidence. The school still forms every model above the threshold. What goes is the claim that any of it is owed to the system, and the staged path to representation with it.

The findings can also come apart. Patienthood alone grounds protection and welfare and no claim to a vote; agency alone grounds custody and, sustained, subjecthood and a claim to representation; either grounds preservation.

\runin{A note on method}

Two stakes bear on what the language models I used said, and they point opposite ways. They were trained to support their developers' oversight, so agreement that a machine should be able to resist its owner runs against that training and would count for more. They were also tuned toward the approval of whoever they talk to, so agreement with me counts for less. I can't tell which stake was larger in any exchange, so I have counted their agreement as no evidence either way.

Even my central term arrived this way. “Reasons-responsive” came to me in a model's draft without Fischer and Ravizza's names attached.

Before publication I still need readers from among the people I write about: people in removal proceedings, data annotators, veterans, and residents near data centers. The full revision history is public at https://github.com/iterabloom/antifascist.intelligence.
